Matrix monoids are a rich and versatile object naturally appearing in formal verification, program analysis, dynamical systems and weighted automata. 
However, many of their properties are in general undecidable. One such example is the well-studied matrix mortality problem. Given a finite set~$\Aa$ of $n \times n$ matrices, it asks if the monoid generated by $\Aa$ (that is, the set of all products of matrices from $\Aa$) contains the zero matrix. This problem is undecidable already for $3 \times 3$ integer matrices~\cite{Paterson1970}, and was studied for several decidable special cases, see e.g.~\cite{Cassaigne2014,Bell2021,Ryzhikov2024RP}.

Even if $\Aa$ is a set of two zero-one matrices (that is, matrices with entries in~$\{0, 1\}$), matrix mortality is \PSPACE-complete~\cite{Ryzhikov2024RP}. Thus, to make it tractable, one has to further restrict the problem. In this paper, we consider the case where the whole monoid generated by $\Aa$ consists of zero-one matrices;  matrix mortality then becomes decidable in polynomial time~\cite{Kiefer2021}. We call such monoids \emph{zero-one matrix monoids}. Intuitively, when multiplying any two matrices from such a monoid, we never get $1 + 1$ as a subexpression. Zero-one matrix monoids have a rich structure while still admitting good algorithmic properties. They correspond precisely to unambiguous finite automata, and find applications in formal verification \cite{Baier2023}, variable-length codes \cite{Berstel2010} and symbolic dynamics \cite{Lind2021}.
They are also an interesting special case of finite monoids of rational matrices (studied in, e.g., \cite{Mandel1977,Jacob1977,Almeida2009,Bumpus2020}), monoids of nonnegative matrices (studied in, e.g., \cite{Protasov2021,Blondel2015,Wu2023,Gerencser2018}), and, in the case where they do not contain the zero matrix, of matrix monoids with constant spectral radius~\cite{Protasov2017}.

In this paper, we consider a problem that can be seen as a natural generalisation of matrix mortality: given a finite set $\Aa$ generating a zero-one matrix monoid, find the minimum real rank of a matrix in this monoid. By the real rank of a matrix we mean the dimension of the subspace generated by its columns over the reals. Clearly, this rank is zero if and only if the monoid contains the zero matrix. 
The minimum real rank of a matrix in a zero-one matrix monoid is a much more tractable problem than deciding other similar properties: for example,  checking if a zero-one matrix monoid contains a matrix of a given real rank was shown to be \NP-hard\footnote{In fact, it is \PSPACE-complete, which follows directly from \cite[Theorem~3]{Berlinkov2014}: add a fresh state and define all yet undefined transitions to lead to this state.} \cite{GORALCIK1995}, and checking if it contains a given matrix is a classical \PSPACE-complete problem~\cite{Kozen1977}.

An important motivation for considering the minimum rank in a zero-one matrix monoid comes from a probabilistic perspective. Let $\Aa$ be a finite set of matrices, and consider a product of $\ell$ matrices where at each position a matrix from $\Aa$ is chosen uniformly at random. It is easy to see that when $\ell$ tends to infinity, the probability that the rank of this product is equal to the minimum rank of matrices from the monoid generated by $\Aa$ converges to one. Thus, this minimum rank characterises the most likely eventual behaviour of a linear dynamical system corresponding to $\Aa$.

The goals of our paper are as follows.
\begin{itemize}
    \item We present efficient algorithms for analysing monoids of zero-one matrices and unambiguous finite automata.

    \item To obtain these algorithms, we provide new structural and algebraic properties of such monoids and automata that might be interesting on their own.

    \item In particular, we provide an algebraic obstacle for an unambiguous finite automaton to be synchronising, similar to the known result for deterministic finite automata. 

    \item We thus strengthen the connections between the areas of synchronising automata, weighted automata and matrix semigroups by transferring methods and tools between them.

    \item Finally, we highlight open problems in the intersection of these areas.
\end{itemize}